\Askhsh{34024_SOLUTION}{1}
\begin{ylh}{1}
    \item [Ύλη από εικόνα]
\end{ylh}
\begin{alist}
\item Λύση
\begin{rlist}
\item Είναι $D_f = \mathbb{R}$ με $f$ συνεχή ως πηλίκο συνεχών συναρτήσεων.
\begin{rlist}
\item Για $x_1, x_2 \in \mathbb{R}$ έχουμε
\[
x_1 < x_2 \Rightarrow e^{x_1} < e^{x_2} \Rightarrow e^{2x_1} < e^{2x_2} \xrightarrow{e^{2x}>0} \frac{1}{e^{2x_1}} > \frac{1}{e^{2x_2}} \Rightarrow f(x_1) > f(x_2) \Rightarrow f \searrow.
\]
Δηλαδή η συνάρτηση $f$ είναι γνησίως φθίνουσα στο $\mathbb{R}$.
\item Έχουμε $\lim_{x \to -\infty} \frac{1}{e^{2x}} = +\infty$ αφού $\lim_{x \to -\infty} e^{2x} = 0$ και $e^{2x} > 0$ κοντά στο $-\infty$,
και $\lim_{x \to +\infty} \frac{1}{e^{2x}} = 0$. Οπότε, λόγω του α)i), παίρνουμε $f(\mathbb{R}) = (0, +\infty)$.
\end{rlist}
\item Από το α)i ερώτημα έχουμε ότι η συνάρτηση $f$ είναι γνησίως μονότονη στο πεδίο ορισμού της, άρα 1-1. Κάθε ευθεία παράλληλη στον άξονα $x'x$ τέμνει την γραφική παράσταση μιας συνάρτησης 1-1 σε ένα το πολύ σημείο. Από το α)ii. ερώτημα η τιμή 3 ανήκει στο σύνολο τιμών της $f$ που είναι μοναδικό, αφού η συνάρτηση είναι 1-1.
\end{alist}